\chapter{The $64^2$ / $64^3$ Globe}

To make operations geometric, the architecture requires a bounded coordinate system that bridges geographical space and operational state.

\section{The $64^2$ Attention Geometry}

A fundamental structural constant in the ecosystem is $64^2 = 4,096$. This provides a practical attention globe, representing the first useful routing surface. It models the world as a bounded operational grid ($\text{region} \times \text{corridor}$).

This $64^2$ grid is not a map tile system in the traditional GIS sense. It is a \textbf{process-control surface}. It geometrically answers critical operational questions: Where is attention required? Which cells are blocked by strikes or storms? Which cells represent lawful capacity? Which cells are candidates for rerouting? 

\section{The $64^3$ Operational Depth}

The breakthrough occurs at $64^3 = 262,144$. The globe transitions from a surface position to operational depth: $\text{region} \times \text{corridor} \times \text{local process place}$. This means a single geographic cell contains up to $64$ local process positions.

\section{GlobeCell(domain, cell, place)}

A route evaluates whether an entity can lawfully move through a coordinate-state trajectory. It answers the geometric question: Can the vessel lawfully move through this sequence of coordinates and states?

\section{Supply-Chain Coordinate States}

In a supply chain context, a vessel is no longer a dot moving across a map. It occupies specific states within a $64^3$ cell. Examples of these coordinate states include:
\begin{itemize}
    \item Port cell + berth state
    \item Port cell + customs state
    \item Port cell + cold-chain state
    \item Port cell + legal state
    \item Port cell + capability state
\end{itemize}
This maps the abstract math directly into tangible supply-chain operations.